\setproblem{21317}

\begin{frame}[label=p2]{징검다리 건너기(\prob)}
\small

\begin{block}{문제}
심마니 영재는 산삼을 찾아다닌다.

산삼을 찾던 영재는 $N$($1 \le N \le 20$)개의 돌이 일렬로 나열되어 있는 강가를 발견했고, 마지막 돌 틈 사이에 산삼이 있다는 사실을 알게 되었다.

마지막 돌 틈 사이에 있는 산삼을 캐기 위해 영재는 돌과 돌 사이를 점프하면서 이동하며, 점프의 종류는 3가지가 있다. 현재 위치에서 다음 돌로 이동하는 작은 점프, 1개의 돌을 건너뛰어 이동하는 큰 점프, 2개의 돌을 건너뛰어 이동하는 매우 큰 점프가 있다. \\

각 점프를 할 때는 에너지를 소비하는데, 작은 점프와 큰 점프 시 소비되는 에너지는 점프를 하는 돌의 번호마다 다르다. 매우 큰 점프는 단 한 번의 기회가 주어지며, 이때는 점프를 하는 돌의 번호와 상관없이 $K$($1 \le K \le 5\,000$)만큼의 에너지를 소비한다. 작은 점프와 큰 점프 시 필요한 에너지도 $5\,000$을 넘지 않는 자연수이다. \\

에너지를 최대한 아껴야 하는 영재가 첫 번째 돌에서 출발하여 산삼을 얻기 위해 필요한 최소 에너지를 구하여라.
\end{block}
\end{frame}

\begin{frame}[fragile]{예제 입력/출력 1}
\begin{columns}[T]
\begin{column}{0.48\textwidth}
\begin{block}{입력}
\begin{verbatim}
5
1 2
2 3
4 5
6 7
4
\end{verbatim}
\end{block}
\end{column}
\begin{column}{0.48\textwidth}
\begin{block}{출력}
\begin{verbatim}
5
\end{verbatim}
\end{block}
\end{column}
\end{columns}
\end{frame}

\begin{frame}[fragile, allowframebreaks]{김시온 소스 코드}
\inputminted{java}{../src/\prob/sion/Main.java}
\end{frame}

\begin{frame}[fragile, allowframebreaks]{원찬혁 소스 코드}
\inputminted{java}{../src/\prob/chanhyeok/Main.java}
\end{frame}

\begin{frame}[fragile, allowframebreaks]{손현준 소스 코드}
\inputminted{java}{../src/\prob/hyunjoon/P21317.java}
\end{frame}

\begin{frame}[fragile, allowframebreaks]{김준형 소스 코드}
\inputminted{java}{../src/\prob/joonhyeong/Main.java}
\end{frame}

\begin{frame}[fragile, allowframebreaks]{서민종 소스 코드}
\inputminted{java}{../src/\prob/minjong/Main.java}
\end{frame}
